This thesis presents the design and implementation of BlockQuiz, a web-based block programming learning platform with
automated grading capabilities, specifically designed for children aged 8-12 years.
As computational thinking becomes increasingly integrated into educational curricula,
notably through Austria's mandatory ``Digitale Grundbildung'' subject introduced in 2022,
there is a growing need for accessible, teacher-friendly programming tools that provide immediate feedback
without requiring extensive technical expertise from educators.

Unlike existing platforms such as Scratch, Tynker, or Microsoft MakeCode, which tend toward
open-ended creative projects or hardware-specific applications, BlockQuiz focuses on short,
tightly constrained programming puzzles with clear learning objectives.
The platform implements two primary exercise types: Input/Output puzzles for algorithmic
thinking and 2D Graphics exercises for visual reasoning and geometry concepts.
Each exercise features a constrained block toolbox, reducing cognitive load while guiding learners toward solutions.

The technical implementation leverages modern web technologies including SvelteKit,
Google Blockly for the visual programming interface, and a privacy-first architecture
enabling self-hosted deployment via Docker. Key contributions include a secure client-side
code execution sandbox with timeout protection, a flexible auto-grading system supporting multiple
validation strategies, and a teacher-oriented content management system for exercise authoring without programming knowledge.

The platform supports internationalization (German and English), implements progressive hint systems,
and captures learning analytics for educational research. An evaluation with target users demonstrates
the platform's usability and effectiveness in teaching core programming concepts including sequences,
loops, conditionals, and simple functions.
